% ===================================================================
%  పట్టిక సంఖ్య 13 — హోటలు. భోజనశాల. కాఫీశాల.
%
%  Traduction, faite en 2026, en telougou d'aujourd'hui, sur l'ido de
%  texto/io. Regles de la colonne : voir 10-pattika-01.tex. Recit a la
%  premiere personne, sans scenes ; les cles ne portent ni c1 ni c2.
%  Les deux notes d'auteur sont en gras-italique dans le fac-simile,
%  contrairement a celles des tableaux 1, 5 et 12 : la colonne suit,
%  comme le hindi.
%
%  DIX-NEUF RETOURNEMENTS SUR TRENTE-TROIS PAIRES — LE RECORD DU
%  LIVRET, et il etait annonce. La colonne marathe avait releve, au
%  meme tableau 13, qu'il coutait neuf inversions pour un tiers du
%  texte du tableau 14, et en avait tire la regle : ce n'est pas la
%  longueur d'un tableau qui coute, c'est sa SYNTAXE. Le tableau 14
%  enumere des metiers ; celui-ci ATTACHE — le conducteur DE
%  l'omnibus, la malle DU voyageur, la cuisine AU sous-sol, l'enseigne
%  DE la halle, l'etendard AU-DESSUS de l'enseigne, la caissiere QUI
%  inscrivait la facture SUR le registre. Chaque « de » et chaque
%  relative est une inversion en telougou.
%
%  Le chiffre le dit mieux que la regle : 53 % de paires retournees au
%  tableau 9, 53 % au 10, 50 % au 11, 46 % au 12, et 58 % ici. Deux
%  langues sans parente — une indo-aryenne, une dravidienne — sont
%  tombees sur le meme tableau comme sur le plus cher des seize, pour
%  la meme raison, et la raison n'est pas dans les langues : elle est
%  dans le texte de 1926.
%
%  LE MENU RESTE FRANCAIS, ET C'EST LA REGLE, PAS UNE PARESSE. Les
%  tripes a la mode de Caen, le gigot de mouton, le Saint-Emilion sont
%  les CHOSES qu'on mange dans cette salle-la ; les rendre par des
%  plats andhra aurait deplace le repas, pas traduit la langue. Le
%  telougou nomme donc ce qu'il voit : కాన్ పద్ధతిలో వండిన పేగులు,
%  గొర్రె తొడ, మలాయితో పాలకూర, రొయ్యలు, వెన్న, జున్ను. Meme choix
%  qu'en hindi, et pour le meme motif.
%
%  LE VOCABULAIRE DE L'HOTEL EST ENTIER EN TELOUGOU, celui du garage
%  ne l'est pas, et la coupure passe la ou on l'attend — a l'annee.
%  L'hotellerie est vieille en Andhra : ద్వారపాలకుడు le portier,
%  వంటవాడు le cuisinier, పనిమనిషి la servante, మోతకూలీ le porteur,
%  బట్టలుతికే అమ్మాయి la blanchisseuse, ఉత్తరాల పెట్టె la boite aux
%  lettres. L'automobile ne l'est pas : డ్రైవరు, మెకానికు, కారు,
%  లిఫ్టు, సైకిలిస్టు, బిలియర్డ్సు s'ecrivent en anglais parce que
%  c'est ainsi qu'on les dit.
%
%  UN JEU QUE JE NE SAIS PAS NOMMER, et je prefere le dire : les
%  « dami » (79), les dames. Le telougou a ses jeux de plateau —
%  వామన గుంటలు, పులిజూదం — mais aucun n'est le jeu de dames, et
%  aucun mot telougou ne le designe. La colonne ecrit చెకర్లు,
%  l'emprunt courant, plutot que de coller le nom d'un jeu indien sur
%  un autre jeu. Meme raison qu'aux trois oiseaux du tableau 8.
% ===================================================================

%%K t13-noto-1 noto
\VUnotes{143.2mm}{%
\textit{\VUgras{(*) పడక గది వివరాల కోసం పట్టిక సంఖ్య 6 చూడండి.}}%
}

%%K t13-apar-1 apar
\VUsaut{3.0mm}
\VUtitre{15.0pt}{60}{పట్టిక సంఖ్య 13}
\VUsaut{1.6mm}
\VUpk{10.2pt}{40}{హోటలు. --- భోజనశాల.}
\VUsaut{2.0mm}
\VUpk{10.2pt}{40}{కాఫీశాల.}
\VUsaut{3.4mm}

%%K t13-01-1 p
1. --- నిన్న సాయంత్రం రోయాన్‌కు చేరుకుని, «\textit{ఓషన్ హోటలు}»లో దిగాను; ఆ హోటలును ఒక
స్నేహితుడు నాకు సిఫారసు చేశాడు.

%%K t13-01-2 p
నాకు అందమైన \VUgras{గది} \textsuperscript{(2)} ఇచ్చారు — రెండో \VUgras{అంతస్తులో}
\textsuperscript{(1)} —, అక్కడ నేను చాలా బాగా నిద్రపోయాను (*).

%%K t13-02-1 p
2. --- ఈ ఉదయం నా గదిలోంచి బయటికి వచ్చాను — అందులోకి \VUgras{పనిమనిషి}
\textsuperscript{(3)} అల్పాహారం తెచ్చింది —, \VUgras{పొగతాగే గదిలోకి} \textsuperscript{(5)}
వెళ్ళాను; అది \VUgras{చదువుకునే గదిగా} \textsuperscript{(4)} కూడా వాడతారు,
\VUgras{పత్రికలు} \textsuperscript{(6)} చదివేందుకు — \VUgras{సిగరెట్టు}
\textsuperscript{(8)} తాగుతూ. చదవడం అయ్యాక \VUgras{లిఫ్టులో} \textsuperscript{(9)} కిందికి
దిగి ఊళ్ళో తిరగడానికి వెళ్ళాను.

%%K t13-03-1 p
3. --- హోటలు \VUgras{గుమాస్తాను} \textsuperscript{(12)} అడగడం మర్చిపోయాను —
\VUgras{ఉమ్మడి భోజన బల్ల} \textsuperscript{(11)} మీద భోజనం ఎన్ని గంటలకు పెడతారో —; అందుకని
త్వరగా తిరిగి వచ్చి, ఆ అవకాశంతో హోటలు \VUgras{స్నానపు గదిలో} \textsuperscript{(10)} స్నానం
చేశాను. తర్వాత నా స్నేహితుడికి ఫోను చేసేందుకు మళ్ళీ \VUgras{కాషు కౌంటరు}
\textsuperscript{(15)} దగ్గరికి వెళ్ళాను. కానీ \VUgras{టెలిఫోను} \textsuperscript{(13)}
పనిలో ఉంది; నా ఇబ్బంది చూసి \VUgras{కాషియరు} \textsuperscript{(14)} — ఆమె అప్పుడు
\VUgras{బిల్లును} \textsuperscript{(17)} \VUgras{ప్రయాణికుల నమోదు పుస్తకంలో}
\textsuperscript{(16)} రాస్తోంది — \VUgras{ద్వారపాలకుడిని} \textsuperscript{(18)} పిలిచి
\VUgras{పనికుర్రాడిని} \textsuperscript{(19)} పంపమని, \VUgras{సైకిలుపై}
\textsuperscript{(20)} నా పని చేయించమని కోరింది.

%%K t13-04-1 p
4. --- ఆమెకు కృతజ్ఞతలు చెప్పి, \VUgras{మోతకూలీకి} \textsuperscript{(24)} (కూలిపని
చేసేవాడు) బహుమానం ఇచ్చేందుకు బయటికి వెళ్ళాను — అతను స్టేషనునుంచి తన \VUgras{చేతిబండిపై}
\textsuperscript{(25)} నా \VUgras{టోపీ డబ్బాను} \textsuperscript{(26)},
\VUgras{నా సూట్‌కేసును} \textsuperscript{(27)}, \VUgras{నా ప్రయాణపు సంచిని}
\textsuperscript{(28)} తెచ్చాడు. ఇప్పుడే వచ్చిన \VUgras{పోస్టుమాన్} \textsuperscript{(21)}
నాకు \VUgras{ఉత్తరం} \textsuperscript{(22)} ఇచ్చాడు.

%%K t13-05-1 p
5. --- \VUgras{ఉత్తరాల పెట్టె} \textsuperscript{(23)} పక్కన, తలుపు గడపమీద నేను ఇంకొంతసేపు
నిలబడి వీధిలో రాకపోకలు చూశాను. \VUgras{నడిపేవాడు} \textsuperscript{(30)} —
\VUgras{గుర్రపు బస్సుది} \textsuperscript{(29)} — తన బండిపై \VUgras{పెద్ద పెట్టెను}
\textsuperscript{(31)} ఎక్కిస్తున్నాడు; అది \VUgras{ప్రయాణికుడిది} \textsuperscript{(32)},
అతన్ని \VUgras{హోటలు యజమాని} \textsuperscript{(33)} సాగనంపుతున్నాడు. కుర్ర
\VUgras{తంతి కుర్రాడు} \textsuperscript{(35)} తీరిగ్గా \VUgras{తంతిని}
\textsuperscript{(34)} హోటలు ఖాతాదారు ఒకరికి తెచ్చాడు; \VUgras{పర్యాటకుడు}
\textsuperscript{(36)} ఒకడు — భుజాన \VUgras{ఫొటో కెమెరాతో} \textsuperscript{(38)} — తన
\VUgras{దారి పుస్తకాన్ని} \textsuperscript{(37)} చూసుకుంటున్నాడు.

%%K t13-tit-1 sub
\VUsaut{4.0mm}
\VUpk{10.2pt}{40}{భోజన వేళ.}
\VUsaut{3.0mm}

%%K t13-06-1 p
6. --- కంచె ముందు, నిర్లక్ష్యంగా \VUgras{పనులు చేసేవాడు} \textsuperscript{(39)} ఒకడు తన
\VUgras{మోత కొక్కేనికీ} \textsuperscript{(40)}, తన \VUgras{పాలిషు పెట్టెకూ}
\textsuperscript{(41)} మధ్య కునికిపాట్లు పడుతున్నాడు; ఈలోగా కుర్ర
\VUgras{బట్టలుతికే అమ్మాయి} \textsuperscript{(42)} — ఆమె బట్టల \VUgras{బుట్ట}
\textsuperscript{(43)} మోసుకెళ్తోంది — తలుపు పక్కన నిలబడిన \VUgras{భోజనశాల పెద్ద}
\textsuperscript{(44)} గంభీరమైన మొహం చూసి నవ్వింది.

%%K t13-07-1 p
7. --- \VUgras{వంటవాడు} \textsuperscript{(45)} కనిపించాడు — అతను తన \VUgras{వంటగదిలోంచి}
\textsuperscript{(47)} బయటికి వచ్చాడు, ఆ వంటగది \VUgras{నేలమాళిగలో} \textsuperscript{(48)}
ఉంది, \VUgras{వంటశాల కుర్రాడిని} \textsuperscript{(46)} పనిమీద పంపేందుకు —; అప్పుడే భోజన
వేళ దగ్గరపడిందని గుర్తొచ్చి భోజనశాలలోకి వెళ్ళాను, దారిలో ఆవరణ దాటుతూ. అక్కడ
\VUgras{డ్రైవరు} \textsuperscript{(50)} తన \VUgras{కారును} \textsuperscript{(49)}
నిలిపాడు; \VUgras{మెకానికు} \textsuperscript{(52)} ఒకడు \VUgras{సైకిలిస్టు}
\textsuperscript{(53)} చెప్పిన ప్రకారం, ఇప్పుడే ప్రమాదానికి గురైన
\VUgras{పెట్రోలు మూడు చక్రాల బండిని} \textsuperscript{(51)} బాగుచేస్తున్నాడు.

%%K t13-08-1 p
8. --- కుర్ర \VUgras{పనికుర్రాడిని} \textsuperscript{(54)} ఒకడిని అడిగాను — అతను
\VUgras{బీరు గ్లాసులు} \textsuperscript{(55)} మోసుకెళ్తున్నాడు — ఉమ్మడి భోజన బల్ల వసారా
కింద ఉందా అని;

%%K t13-08-1 p suite
అతను అవునన్నందుకు, ఒక బల్ల దగ్గర చోటు తీసుకుని చక్కటి భోజనం చేయించుకున్నాను.

%%K t13-noto-2 noto
\VUnotes{143.2mm}{%
\textit{\VUgras{(*) పదకోశంలో ఇచ్చిన వంటకాల జాబితా సాయంతో ఉపాధ్యాయుడు కింది
భోజన పత్రాన్ని సులభంగా మార్చుకోవచ్చు.}}%
}

%%K t13-tit-2 sub
\VUsaut{4.0mm}
\VUpk{10.2pt}{40}{చక్కటి భోజనం.}
\VUsaut{3.0mm}

%%K t13-09-1 p
9. --- సర్వరు నాకు మధ్యాహ్న భోజన \VUgras{పత్రాన్ని} (*), ద్రాక్షసారా జాబితానూ తెచ్చాడు.
నేను ఎంచుకుని, \VUgras{ముందు వడ్డనగా} \VUgras{వెన్నతో} \VUgras{రొయ్యలు}, మొదటి
వడ్డనగా \VUgras{కాన్ పద్ధతిలో వండిన పేగులు} ఆర్డరు చేశాను. \VUgras{చేప} తినలేదు, కానీ
గొర్రె \VUgras{తొడ} ముక్క అడిగాను, \VUgras{కూరగా} మలాయితో \VUgras{పాలకూర}.
ఆ తర్వాత \VUgras{మధ్య వంటకాల్లో} ఏదీ నచ్చకపోవడంతో రకరకాల తీపి వంటకాలు తెప్పించుకున్నాను :
\VUgras{జున్ను}, \VUgras{పండ్లు}, \VUgras{బిస్కెట్లు}. వాటికి చక్కటి సెంతెమిలియోఁ
\VUgras{వైన్} జోడించాను; భోజనంతో బాగా సంతృప్తి చెంది, \VUgras{సర్వరుకు}
\textsuperscript{(57)} మంచి \VUgras{టిప్పు} \textsuperscript{(59)} ఇచ్చి బల్ల వదిలాను.

%%K t13-10-1 p
10. --- నేను బయలుదేరేందుకు సిద్ధంగా ఉండటం చూసి, \VUgras{నిర్వాహకుడు}
\textsuperscript{(60)} బాల్కనీలోకి వెళ్ళమని సూచించాడు — \VUgras{మహాసముద్రం}
\textsuperscript{(62)} అద్భుత దృశ్యం చూసేందుకు. దూరంగా \VUgras{చేపల చిన్న పడవలు}
\textsuperscript{(63)}, \VUgras{తెరచాప ఓడలు} \textsuperscript{(64)}, భారీ
\VUgras{ప్రయాణ ఓడ} \textsuperscript{(65)} కనిపించడం మొదలైంది — ఆ ఓడ నల్లని నీడ తెల్లని
\VUgras{మేఘాలపై} \textsuperscript{(67)} పొడుచుకొచ్చి కనిపిస్తుంది, \VUgras{క్షితిజపు}
\textsuperscript{(66)} మేఘాలవి. తేలికపాటి గాలి \VUgras{పతాకాన్ని} \textsuperscript{(70)}
ఊపింది — అది \VUgras{బోర్డు} \textsuperscript{(69)} మీద నాటినది,
\VUgras{సంత మంటపపు} \textsuperscript{(68)} బోర్డు.

%%K t13-tit-3 sub
\VUsaut{3.0mm}
\VUpk{10.2pt}{40}{కాలక్షేపాలు.}
\VUsaut{3.0mm}

%%K t13-11-1 p
11. --- ఆ దృశ్యం ఎంత ఆసక్తికరంగా ఉందంటే, రేపు \VUgras{బైనాక్యులర్లు}
\textsuperscript{(71)} గానీ \VUgras{దూరదర్శిని} \textsuperscript{(72)} గానీ పట్టుకుని
కాఫీశాల మిద్దె మీదికి ఎక్కాలని నిర్ణయించుకున్నాను.

%%K t13-12-1 p
12. --- బాల్కనీ వదిలి, హోటలు కాఫీశాలకు వెళ్ళాను — \VUgras{పానీయం}
\textsuperscript{(73)} తాగేందుకు, \VUgras{బిలియర్డ్సు ఆట} \textsuperscript{(75)}
ఆడేందుకు. ఆగకుండా కాఫీ గదిని దాటాను; అందులో చాలామంది \VUgras{చదరంగం}
\textsuperscript{(78)}, \VUgras{చెకర్లు} \textsuperscript{(79)},
\VUgras{డొమినో పలకలు} \textsuperscript{(80)}, \VUgras{ట్రిక్‌ట్రాక్}
\textsuperscript{(81)}, \VUgras{పేకాట} \textsuperscript{(82)} ఆడుతున్నారు; నేను నేరుగా
\VUgras{బిలియర్డ్సు గదికి} \textsuperscript{(74)} ఎక్కాను.

%%K t13-13-1 p
13. --- ఆట ముగిశాక కింది అంతస్తుకు దిగి, ఉత్తరం రాసేందుకు సర్వరును \VUgras{కవరు}
\textsuperscript{(84)}, \VUgras{కాగితం} \textsuperscript{(83)},
\VUgras{తపాలా బిళ్ళ} \textsuperscript{(85)}, \VUgras{కలం} \textsuperscript{(87)},
\VUgras{సిరా} \textsuperscript{(86)} అడిగాను.

%%K t13-13-2 p
ఆ తర్వాత \VUgras{బండి తోలేవాడిని} \textsuperscript{(92)} పిలిచాను — అతని \VUgras{బండి}
\textsuperscript{(88)} కాఫీశాల ముందు ఆగి ఉంది. గంటల లెక్కనా, ఒక్క ప్రయాణానికా ధర ఎంతని
అడిగాను; ధర కుదిరాక అతను నాకు \VUgras{బండి తలుపు} \textsuperscript{(89)} తెరిచి
కూర్చోబెట్టాడు. మళ్ళీ \VUgras{తోలుసీటుపైకి} \textsuperscript{(91)} ఎక్కి,
\VUgras{కొరడాతో} \textsuperscript{(93)} గుర్రాలను వేగంగా కదిలించాడు; మేము ముచ్చటైన
విహారానికి బయలుదేరాం.

\VUsaut{6.0mm}
\VUfilet{22mm}
